%% =====================================================================
%%  IEEE SPMB 2026 — Supplement to v4
%%  Appendix A: Full proofs for the ESF Recoverability Framework
%%  (Definition 1, Propositions 1 and 2, Theorem 1)
%% =====================================================================

\documentclass[conference]{IEEEtran}

\IEEEoverridecommandlockouts
\usepackage{amsmath,graphicx,soul,subcaption,balance}
\usepackage[numbers,sort&compress]{natbib}
\renewcommand\bibsection{\section*{\normalsize References}}
\usepackage{float}
\usepackage{times}
\usepackage{multirow}
\usepackage{color}
\usepackage{soul,mathrsfs}
\usepackage[utf8]{inputenc}
\usepackage[T1]{fontenc}
\usepackage{array}
\usepackage[hyphens, spaces, obeyspaces]{url}
\def\UrlFont{\em}
\usepackage{booktabs}
\usepackage{xcolor}
\usepackage{amssymb,amsfonts,amsthm}


\usepackage{caption}
\captionsetup[figure]{font=small,name=Fig.}
\captionsetup[table]{font=small,name=TABLE}

\captionsetup{justification=justified,singlelinecheck=true}
\captionsetup{belowskip=0pt}



\newtheorem{definition}{Definition}
\newtheorem{proposition}{Proposition}
\newtheorem{theorem}{Theorem}
\newtheorem{lemma}{Lemma}

\title{Supplement to ``ESF: Harmonizing Heterogeneous EEG Acquisition for
Pretrained Classifiers'' --- Appendix A: Proofs}
\author{\IEEEauthorblockN{Hitesh Dammu}
\IEEEauthorblockA{Independent Researcher, Hyderabad, India \\ hdammu@asu.edu}}

\renewcommand{\thetable}{S\arabic{table}}
\renewcommand{\thefigure}{S\arabic{figure}}
\newcommand{\PINSKERNUMBERS}{Over the $40$ measured non-identity arms the class-conditional bound holds on $40/40$ for EEGPT and $40/40$ for BIOT\PINSKERLABRAM, and is vacuous (sum of class-conditional TVs above $1$) on the strongly shifted naive arms; the marginal ratio has median $0.15$ (EEGPT) and $0.24$ (BIOT) and maximum $1.41$ and $1.35$ respectively, both on the $128$\,Hz naive arm, so no marginal constant below $1$ exists for any backbone in this sweep.}
\newcommand{\PINSKERLABRAM}{}
\newcommand{\LABRAMABS}{}
\newcommand{\LABRAMDISC}{}
\newcommand{\LABRAMAONE}{}
\newcommand{\LABRAMATWO}{}
\newcommand{\LABRAMATHREE}{}
\newcommand{\LABRAMAFOUR}{}
\newcommand{\LABRAMAFIVE}{}
\newcommand{\BANDWIDTHOTHERS}{}
\newcommand{\ENVELOPETEXT}{}
\newcommand{\LABRAMBASE}{}
\newcommand{\EEGPTPRETRAIN}{(PhysioNet motor imagery, the Tsinghua SSVEP benchmark, M3CV and SEED, per the authors' repository)}
\begin{document}
\maketitle

\section*{Appendix A: Proofs}

\subsection*{A.1 Standing assumptions}
\begin{itemize}\itemsep0pt
\item[(S1)] Signals $s$ are finite-energy real-valued multichannel
functions on $[0,T] \times \{1,\dots,C\}$, identified with elements of
$L^2([0,T])^C$.
\item[(S2)] The canonical band is $\mathcal{B} = [0, f_{\max}]$ with
$f_{\max} \le f_0/2$; the canonical rate is $f_0{=}250$\,Hz, with
$f_{\max}{=}70$\,Hz.
\item[(S3)] The acquisition operator $A_\theta$ is measurable from
$\mathcal{S}$ to a discrete-time space $\hat{\mathcal{S}}_\theta$;
$K_\theta$ is measurable in the other direction.
\item[(S4)] The classifier $f : \mathcal{S} \to \{0,1\}$ is the
composition of a frozen feature map $g : \mathcal{S} \to \mathbb{R}^d$
and a linear probe $w : \mathbb{R}^d \to \mathbb{R}$ thresholded at
$0$. AUROC $\mathcal{A}(f; \mathcal{D})$ is the probability that the
probe ranks a positive example above a negative example,
$\Pr(w(g(X^+)) > w(g(X^-)))$.
\end{itemize}

\subsection*{A.2 Proof of Proposition~1 (left-invertibility $\Rightarrow$ recoverability)}
By assumption $A_\theta$ is left-invertible with inverse
$A_\theta^{-L}$. Set $K_\theta = A_\theta^{-L}$. Then for every
$s \in \mathcal{S}$
\[ K_\theta(A_\theta(s)) = s, \]
hence $f \circ K_\theta \circ A_\theta = f$ pointwise, and the AUROC
functional is unchanged. The gap in Definition~1 is exactly $0$,
i.e., the axis is $(f,0)$-recoverable. $\square$

\paragraph{Specialization to per-channel gain.} $A_g s = g \cdot s$
with $g \in (0, \infty)^C$ is left-invertible by $A_g^{-L} h = g^{-1}
\cdot h$ (channel-wise). Numerical implementation has roundoff
$O(2^{-\text{mantissa}})$, which explains the small residual gap
observed at $g \in \{0.5, 0.75\}$.

\paragraph{Specialization to power-line interference.} The
deterministic sinusoidal injection $A_{\nu,a} s = s + a \sin(2\pi \nu
t)$ is left-invertible by subtracting the sinusoid (or, in the
IIR-notch implementation used here, by a narrowband notch at $\nu$ with
passband leakage bounded by a target $\epsilon_{\text{pass}}$). With
$\Delta \nu = 1$\,Hz and $\nu = 50$\,Hz the passband loss inside
$\mathcal{B}$ is empirically $<10^{-3}$ in the AUROC functional, well
within the $\varepsilon = 0$ regime up to numerical precision.

\subsection*{A.3 Proof of Proposition~2 (sub-Nyquist irrecoverability)}
\label{app:proof_nyquist}
Let $S(\omega) = \int s(t) e^{-i\omega t}\,dt$ denote the Fourier
transform of $s$ (component-wise). The acquisition operator $A_\theta$
consists of (i) the ideal low-pass projection $\Pi_{f_n/2}$ onto
frequencies in $[0, f_n/2]$, and (ii) uniform decimation to $f_n$. By
the Shannon sampling theorem, the decimated representation
$A_\theta(s)$ depends only on $\Pi_{f_n/2}(s)$, that is,
$A_\theta(s) = A_\theta(\Pi_{f_n/2} s)$. Therefore for every
canonicalization $K_\theta$, the composition
$(K_\theta \circ A_\theta)(s) = (K_\theta \circ A_\theta)(\Pi_{f_n/2}
s)$; the image lies in the bandlimited subspace
$\mathcal{S}_{f_n/2} = \{s : \mathrm{supp}(S) \subseteq [0, f_n/2]\}$
regardless of $K_\theta$.

Now consider two test signals $s, s' \in \mathcal{S}$ that agree on
$\Pi_{f_n/2}$ but for which $f(s) \ne f(s')$ (such a pair exists by
hypothesis). Their composition images are identical:
$(f \circ K_\theta \circ A_\theta)(s) = (f \circ K_\theta \circ
A_\theta)(s')$, yet the truth labels (and the clean predictions)
differ. This produces a mandatory label-disagreement contribution to
the AUROC under any $K_\theta$. The AUROC functional restricted to
$\mathcal{S}_{f_n/2}$ is a lower bound for any $K_\theta$, hence
\[
\mathcal{A}(f \circ K_\theta \circ A_\theta; \mathcal{D}) \le
\mathcal{A}(f|_{\mathcal{S}_{f_n/2}}; \mathcal{D}),
\]
and the gap to $\mathcal{A}(f; \mathcal{D})$ is irrecoverable.
$\square$

\paragraph{Empirical status of Proposition~2 on this classifier (v2).}
Proposition~2 is conditional: it applies when $f$ depends on spectral
content in $(f_n/2, f_{\max}]$. The v2 sweep, in which a $128$\,Hz
acquisition is simulated by anti-aliased decimation of the source
recording and then canonicalized back to $250$\,Hz, measures a
canonicalized AUROC of $0.903$ against a baseline of $0.904$
(gap $0.0004$, not significant). The band $(64, 125]$\,Hz therefore
carries no information that this frozen EEGPT probe uses on TUAB, and
the hypothesis of Proposition~2 is not met: the proposition predicts a
nonzero gap only when the lost band is load-bearing, and for this
classifier it is not. The v1 report of a $0.150$ loss at $128$\,Hz was
an artefact of the v1 implementation, which time-warped the
$250$\,Hz signal rather than simulating acquisition (Version note).
The same holds at $f_n{=}512$\,Hz (canonicalized AUROC $0.904$); the
v1 ``on-device anti-alias kernel mismatch'' explanation described a
filter the author had written, not a device property, and is
withdrawn. What the $128$ and $512$\,Hz cells do show is the cost of
\emph{not} canonicalizing: the naive arm, which presents native samples
to the model as if they were at $250$\,Hz, falls to AUROC $0.735$ and
$0.723$ respectively (both FDR-significant), so the sampling-rate axis
is recoverable in practice and the recovery is large.

\paragraph{Extension to broadband Gaussian noise.} Additive
zero-mean broadband Gaussian noise $A_\sigma s = s + n$, $n \sim
\mathcal{N}(0, \sigma^2 I)$, is not left-invertible: from a single
realization $\hat{s} = s + n$ the underlying $s$ is recoverable only
in distribution, not pointwise. Hence no $K$ in our canonicalization
toolbox achieves $K \circ A = \mathrm{id}$, and the conclusion of
Proposition~2 transfers with the lower-bound argument replaced by the
noise-rate--distortion lower bound.

\subsection*{A.4 Proof of Theorem~1 (AUROC--TV bound), corrected statement}
\label{app:proof_tv}
\emph{Correction in v2.} The v1 proof passed from class-conditional total
variation to the marginal total variation in its last step. That step
does not hold in general: the marginal is the mixture
$\pi_+ P|_{+} + \pi_- P|_{-}$, and a shift that moves the two class
conditionals in opposite directions can leave the marginal almost
unchanged while changing the AUROC substantially. The theorem is
therefore restated in class-conditional form; the marginal quantity is
retained only as a monitoring heuristic (below), without a guarantee.

Let $Y \in \{0,1\}$ be the binary label and $Z = w(g(X))$ the bounded
probe score under input $X \sim \mathcal{D}$. Write $P_0|_{y}$ and
$P_\theta|_{y}$ for the class-conditional score distributions under the
clean and the acquisition-shifted input distributions, $y \in \{+,-\}$.
Assume labels are unaffected by acquisition (no label shift). The AUROC
is the probability that an independent positive outranks a negative,
\[
\mathcal{A}(f; \mathcal{D}) = \Pr_{(Z^+, Z^-) \sim P|_{+} \otimes
P|_{-}} [Z^+ > Z^-].
\]

\textbf{Theorem 1 (restated).} For any bounded-score classifier $f$,
writing $\Delta\mathcal{A} = \mathcal{A}(f;\mathcal{D}_0)-\mathcal{A}(f;\mathcal{D}_\theta)$
and $\mathrm{TV}_y = \mathrm{TV}(P_0|_{y},P_\theta|_{y})$,
$D_y = D_{\mathrm{KL}}(P_\theta|_{y}\,\|\,P_0|_{y})$ for $y\in\{+,-\}$,
\begin{align*}
|\Delta\mathcal{A}| &\le \mathrm{TV}_{+} + \mathrm{TV}_{-} \\
&\le \sqrt{\tfrac12 D_{+}} + \sqrt{\tfrac12 D_{-}} .
\end{align*}

\textbf{Proof.} By Strassen's theorem there exist couplings
$(Z^+, Z'^{+})$ with marginals $(P_0|_{+}, P_\theta|_{+})$ and
$(Z^-, Z'^{-})$ with marginals $(P_0|_{-}, P_\theta|_{-})$ such that
$\Pr(Z^+\neq Z'^{+}) = \mathrm{TV}(P_0|_{+},P_\theta|_{+})$ and
$\Pr(Z^-\neq Z'^{-}) = \mathrm{TV}(P_0|_{-},P_\theta|_{-})$; take the two
couplings independent of each other. On the event
$\{Z^+ = Z'^{+}\}\cap\{Z^- = Z'^{-}\}$ the indicators
$\mathbf 1[Z^+>Z^-]$ and $\mathbf 1[Z'^{+}>Z'^{-}]$ coincide, so
\begin{align*}
&\bigl|\Pr[Z^+>Z^-]-\Pr[Z'^{+}>Z'^{-}]\bigr| \\
&\quad\le \Pr\bigl(Z^+\neq Z'^{+}\bigr) + \Pr\bigl(Z^-\neq Z'^{-}\bigr)
 = \mathrm{TV}_{+} + \mathrm{TV}_{-}
\end{align*}
by the union bound. The second inequality is Pinsker applied to each
class conditional~\cite{pinsker1964information, cover2006elements}.
$\square$

\textbf{Remarks.} (i) The constant is $1$ per class, not $2$; the v1
factor $4$ arose from a loose union bound followed by the invalid
mixture step. (ii) The bound is in terms of class-conditional score
distributions and therefore requires labels to evaluate. At deployment
labels are unavailable; the quantity that can be monitored is the
marginal score distribution $P_\theta$. Its distance from $P_0$ is a
\emph{heuristic proxy} for the class-conditional distances: small
marginal shift does not imply small AUROC change, because opposite
class-conditional movements can cancel in the mixture.

\paragraph{Empirical check (Section IV-C).} On the v2 sweep the
class-conditional bound of Theorem~1 holds on all $40$ measured
non-identity arms, as it must; it is loose (the class-conditional total
variations sum to more than $1$ on the strongly shifted arms, where the
bound is vacuous). The marginal ratio $C_{\mathrm{arm}} =
|\Delta\mathrm{AUROC}| / \mathrm{TV}(P_0,P_\theta)$ has median $0.15$
and maximum $1.41$ (the $128$\,Hz naive arm, where a marginal shift of
$0.17$ accompanies an AUROC loss of $0.17$). A marginal constant below
$1$, as v1 reported ($0.47$), does not exist for this sweep; this is the
expected behaviour of the marginal heuristic and the reason the theorem
is stated class-conditionally.

\subsection*{A.5 Composability of axes}
The five-axis sweep treats each axis independently. In deployment, two
or more axes co-occur. Composability follows from operator
composition, $A_{\theta_1, \theta_2} = A_{\theta_2} \circ A_{\theta_1}$,
and recoverability composes under invertibility.
\begin{lemma}\label{lem:compose}
If $A_{\theta_1}$ and $A_{\theta_2}$ are both $(f,0)$-recoverable with
canonicalizations $K_{\theta_1}, K_{\theta_2}$, then their composition
is $(f, 0)$-recoverable with $K = K_{\theta_1} \circ K_{\theta_2}$. If
either component is irrecoverable in the sense of Proposition~2, the
composition is irrecoverable with the same lower-bound argument applied
to the dominant operator.
\end{lemma}

\paragraph{Practical implication.} The lemma orders the deployment
response: a recording whose broadband SNR is below the measured
envelope is not recovered by any combination of the invertible stages
(notch, rescaling, resampling), so an intake gate on the SNR axis is
justified irrespective of the metadata on the other axes, whereas a
recording that is off-contract only on invertible axes is admitted and
canonicalized. In v2 the sub-Nyquist example used in v1 no longer
applies to this classifier, since $128$\,Hz acquisition is inside the
measured envelope (\S IV-B); the composition rule itself is unchanged.

\subsection*{A.6 Relation to prior frameworks}
The Ben-David et al.\ domain-adaptation bound~\cite{bendavid2010theory}
\[
\varepsilon_T(h) \le \varepsilon_S(h) +
\tfrac{1}{2} d_{\mathcal{H}\Delta\mathcal{H}}(\mathcal{D}_S, \mathcal{D}_T) +
\lambda
\]
gives a divergence bound on \emph{target risk} for an arbitrary
classifier in $\mathcal{H}$. Our framework is complementary in three
ways: (i)~the bound is on the \emph{deployed} classifier $f$, not the
worst $h$; (ii)~the divergence is on the per-recording \emph{score}
distribution, computable from a single deployment without unlabelled
target samples; (iii)~the recoverability predicate is a
\emph{structural} property of the acquisition operator that allows us
to predict the partition before collecting any data.


\section*{B. Per-cell metrics for every backbone}
\subsection*{B.1 Per-cell metrics}
Tables~\ref{tab:s1-eegpt}--\ref{tab:s1-biot} give, for each backbone,
the AUROC, balanced accuracy and expected calibration error of both
arms in every cell of the sweep, averaged over the three MLP-probe
seeds. The broadband (A5) and low-pass (B) axes have no corrective
stage, so their two arms coincide by construction. Values are
traceable to \texttt{results/sweep\_<backbone>\_mlp\_s<seed>.csv} in
the released repository.

\begin{table}[htbp]
\caption{EEGPT: per-cell metrics, mean over MLP seeds 7/11/13.}
\label{tab:s1-eegpt}
\centering\footnotesize\setlength{\tabcolsep}{3pt}
\IfFileExists{../agg/table_S1_eegpt.tex}{\begin{tabular}{lcccccc}
\toprule
Condition & \multicolumn{3}{c}{Naive} & \multicolumn{3}{c}{Canonicalized} \\
& AUROC & bal.\ acc. & ECE & AUROC & bal.\ acc. & ECE \\
\midrule
\multicolumn{7}{l}{\emph{(A1) Sampling rate $f_n$ (Hz)}}\\
\quad 128 & 0.729 & 0.613 & 0.228 & 0.900 & 0.806 & 0.096 \\
\quad 200 & 0.879 & 0.782 & 0.115 & 0.900 & 0.816 & 0.087 \\
\quad 256 & 0.892 & 0.791 & 0.117 & 0.900 & 0.817 & 0.082 \\
\quad 512 & 0.727 & 0.565 & 0.373 & 0.900 & 0.818 & 0.081 \\
\midrule
\multicolumn{7}{l}{\emph{(A2) Gain $g$}}\\
\quad 0.5 & 0.820 & 0.658 & 0.312 & 0.900 & 0.817 & 0.082 \\
\quad 0.75 & 0.856 & 0.737 & 0.179 & 0.900 & 0.817 & 0.082 \\
\quad 1.5 & 0.884 & 0.773 & 0.127 & 0.900 & 0.817 & 0.082 \\
\quad 2 & 0.878 & 0.737 & 0.210 & 0.900 & 0.817 & 0.082 \\
\midrule
\multicolumn{7}{l}{\emph{(A3) Mains $a$ ($\mu$V)}}\\
\quad 5 & 0.878 & 0.610 & 0.388 & 0.900 & 0.817 & 0.082 \\
\quad 15 & 0.801 & 0.502 & 0.528 & 0.901 & 0.817 & 0.082 \\
\quad 30 & 0.683 & 0.500 & 0.542 & 0.900 & 0.817 & 0.081 \\
\quad 60 & 0.598 & 0.500 & 0.543 & 0.900 & 0.817 & 0.081 \\
\midrule
\multicolumn{7}{l}{\emph{(A4) Montage}}\\
\quad legacy16 & 0.893 & 0.762 & 0.174 & 0.888 & 0.767 & 0.142 \\
\quad bipolar18 & 0.878 & 0.700 & 0.247 & 0.880 & 0.701 & 0.241 \\
\midrule
\multicolumn{7}{l}{\emph{(A5) Broadband SNR (dB)}}\\
\quad -5 & 0.737 & 0.507 & 0.355 & 0.737 & 0.507 & 0.355 \\
\quad 0 & 0.792 & 0.547 & 0.327 & 0.792 & 0.547 & 0.327 \\
\quad 5 & 0.841 & 0.631 & 0.263 & 0.841 & 0.631 & 0.263 \\
\quad 10 & 0.863 & 0.694 & 0.191 & 0.863 & 0.694 & 0.191 \\
\quad 20 & 0.890 & 0.778 & 0.108 & 0.890 & 0.778 & 0.108 \\
\midrule
\multicolumn{7}{l}{\emph{(B) Low-pass cutoff (Hz)}}\\
\quad 15 & -- & -- & -- & 0.900 & 0.654 & 0.311 \\
\quad 20 & -- & -- & -- & 0.899 & 0.776 & 0.163 \\
\quad 25 & -- & -- & -- & 0.901 & 0.810 & 0.099 \\
\quad 30 & -- & -- & -- & 0.902 & 0.814 & 0.074 \\
\quad 40 & -- & -- & -- & 0.902 & 0.820 & 0.073 \\
\quad 50 & -- & -- & -- & 0.903 & 0.817 & 0.073 \\
\quad 64 & -- & -- & -- & 0.899 & 0.806 & 0.093 \\
\quad 80 & -- & -- & -- & 0.899 & 0.803 & 0.096 \\
\quad 100 & -- & -- & -- & 0.900 & 0.818 & 0.084 \\
\bottomrule
\end{tabular}
}{\emph{pending}}
\end{table}

\begin{table}[htbp]
\caption{LaBraM: per-cell metrics, mean over MLP seeds 7/11/13.}
\label{tab:s1-labram}
\centering\footnotesize\setlength{\tabcolsep}{3pt}
\IfFileExists{../agg/table_S1_labram.tex}{\begin{tabular}{lcccccc}
\toprule
Condition & \multicolumn{3}{c}{Naive} & \multicolumn{3}{c}{Canonicalized} \\
& AUROC & bal.\ acc. & ECE & AUROC & bal.\ acc. & ECE \\
\midrule
\multicolumn{7}{l}{\emph{(A1) Sampling rate $f_n$ (Hz)}}\\
\quad 128 & 0.813 & 0.673 & 0.174 & 0.815 & 0.503 & 0.473 \\
\quad 200 & 0.845 & 0.768 & 0.063 & 0.898 & 0.812 & 0.071 \\
\quad 256 & 0.889 & 0.792 & 0.137 & 0.895 & 0.800 & 0.111 \\
\quad 512 & 0.608 & 0.492 & 0.483 & 0.895 & 0.795 & 0.110 \\
\midrule
\multicolumn{7}{l}{\emph{(A2) Gain $g$}}\\
\quad 0.5 & 0.875 & 0.786 & 0.084 & 0.894 & 0.796 & 0.111 \\
\quad 0.75 & 0.878 & 0.787 & 0.093 & 0.894 & 0.796 & 0.111 \\
\quad 1.5 & 0.884 & 0.795 & 0.097 & 0.894 & 0.796 & 0.111 \\
\quad 2 & 0.884 & 0.793 & 0.100 & 0.894 & 0.796 & 0.111 \\
\midrule
\multicolumn{7}{l}{\emph{(A3) Mains $a$ ($\mu$V)}}\\
\quad 5 & 0.863 & 0.782 & 0.075 & 0.895 & 0.796 & 0.111 \\
\quad 15 & 0.801 & 0.688 & 0.211 & 0.895 & 0.796 & 0.111 \\
\quad 30 & 0.786 & 0.585 & 0.331 & 0.894 & 0.798 & 0.111 \\
\quad 60 & 0.681 & 0.532 & 0.389 & 0.894 & 0.797 & 0.110 \\
\midrule
\multicolumn{7}{l}{\emph{(A4) Montage}}\\
\quad legacy16 & 0.851 & 0.530 & 0.412 & 0.850 & 0.534 & 0.403 \\
\quad bipolar18 & 0.874 & 0.643 & 0.271 & 0.875 & 0.643 & 0.274 \\
\midrule
\multicolumn{7}{l}{\emph{(A5) Broadband SNR (dB)}}\\
\quad -5 & 0.691 & 0.500 & 0.477 & 0.691 & 0.500 & 0.477 \\
\quad 0 & 0.744 & 0.500 & 0.465 & 0.744 & 0.500 & 0.465 \\
\quad 5 & 0.777 & 0.511 & 0.432 & 0.777 & 0.511 & 0.432 \\
\quad 10 & 0.809 & 0.566 & 0.373 & 0.809 & 0.566 & 0.373 \\
\quad 20 & 0.847 & 0.686 & 0.246 & 0.847 & 0.686 & 0.246 \\
\midrule
\multicolumn{7}{l}{\emph{(B) Low-pass cutoff (Hz)}}\\
\quad 15 & 0.864 & 0.500 & 0.542 & 0.864 & 0.500 & 0.542 \\
\quad 20 & 0.845 & 0.500 & 0.542 & 0.845 & 0.500 & 0.542 \\
\quad 25 & 0.840 & 0.500 & 0.540 & 0.840 & 0.500 & 0.540 \\
\quad 30 & 0.838 & 0.500 & 0.533 & 0.838 & 0.500 & 0.533 \\
\quad 40 & 0.827 & 0.496 & 0.516 & 0.827 & 0.496 & 0.516 \\
\quad 50 & 0.821 & 0.496 & 0.504 & 0.821 & 0.496 & 0.504 \\
\quad 64 & 0.815 & 0.503 & 0.470 & 0.815 & 0.503 & 0.470 \\
\quad 80 & 0.850 & 0.555 & 0.376 & 0.850 & 0.555 & 0.376 \\
\quad 100 & 0.894 & 0.806 & 0.076 & 0.894 & 0.806 & 0.076 \\
\bottomrule
\end{tabular}
}{\emph{pending}}
\end{table}

\begin{table}[htbp]
\caption{BIOT: per-cell metrics, mean over MLP seeds 7/11/13.}
\label{tab:s1-biot}
\centering\footnotesize\setlength{\tabcolsep}{3pt}
\IfFileExists{../agg/table_S1_biot.tex}{\begin{tabular}{lcccccc}
\toprule
Condition & \multicolumn{3}{c}{Naive} & \multicolumn{3}{c}{Canonicalized} \\
& AUROC & bal.\ acc. & ECE & AUROC & bal.\ acc. & ECE \\
\midrule
\multicolumn{7}{l}{\emph{(A1) Sampling rate $f_n$ (Hz)}}\\
\quad 128 & 0.575 & 0.560 & 0.324 & 0.668 & 0.501 & 0.494 \\
\quad 200 & 0.845 & 0.776 & 0.086 & 0.893 & 0.807 & 0.092 \\
\quad 256 & 0.896 & 0.826 & 0.107 & 0.898 & 0.826 & 0.060 \\
\quad 512 & 0.403 & 0.500 & 0.540 & 0.898 & 0.827 & 0.061 \\
\midrule
\multicolumn{7}{l}{\emph{(A2) Gain $g$}}\\
\quad 0.5 & 0.836 & 0.618 & 0.317 & 0.898 & 0.827 & 0.062 \\
\quad 0.75 & 0.845 & 0.708 & 0.189 & 0.898 & 0.827 & 0.062 \\
\quad 1.5 & 0.875 & 0.775 & 0.092 & 0.898 & 0.827 & 0.062 \\
\quad 2 & 0.857 & 0.749 & 0.173 & 0.898 & 0.827 & 0.062 \\
\midrule
\multicolumn{7}{l}{\emph{(A3) Mains $a$ ($\mu$V)}}\\
\quad 5 & 0.661 & 0.500 & 0.542 & 0.898 & 0.827 & 0.061 \\
\quad 15 & 0.633 & 0.500 & 0.543 & 0.898 & 0.827 & 0.061 \\
\quad 30 & 0.643 & 0.500 & 0.543 & 0.898 & 0.827 & 0.062 \\
\quad 60 & 0.627 & 0.500 & 0.543 & 0.898 & 0.826 & 0.063 \\
\midrule
\multicolumn{7}{l}{\emph{(A4) Montage}}\\
\quad legacy16 & 0.873 & 0.577 & 0.361 & 0.867 & 0.584 & 0.357 \\
\quad bipolar18 & 0.878 & 0.796 & 0.066 & 0.879 & 0.795 & 0.073 \\
\midrule
\multicolumn{7}{l}{\emph{(A5) Broadband SNR (dB)}}\\
\quad -5 & 0.678 & 0.509 & 0.280 & 0.678 & 0.509 & 0.280 \\
\quad 0 & 0.785 & 0.529 & 0.290 & 0.785 & 0.529 & 0.290 \\
\quad 5 & 0.810 & 0.571 & 0.267 & 0.810 & 0.571 & 0.267 \\
\quad 10 & 0.842 & 0.620 & 0.246 & 0.842 & 0.620 & 0.246 \\
\quad 20 & 0.847 & 0.717 & 0.133 & 0.847 & 0.717 & 0.133 \\
\midrule
\multicolumn{7}{l}{\emph{(B) Low-pass cutoff (Hz)}}\\
\quad 15 & -- & -- & -- & 0.814 & 0.603 & 0.335 \\
\quad 20 & -- & -- & -- & 0.656 & 0.500 & 0.542 \\
\quad 25 & -- & -- & -- & 0.820 & 0.500 & 0.539 \\
\quad 30 & -- & -- & -- & 0.844 & 0.512 & 0.460 \\
\quad 40 & -- & -- & -- & 0.813 & 0.611 & 0.314 \\
\quad 50 & -- & -- & -- & 0.702 & 0.503 & 0.503 \\
\quad 64 & -- & -- & -- & 0.681 & 0.505 & 0.486 \\
\quad 80 & -- & -- & -- & 0.807 & 0.605 & 0.374 \\
\quad 100 & -- & -- & -- & 0.897 & 0.820 & 0.066 \\
\bottomrule
\end{tabular}
}{\emph{pending}}
\end{table}


\section*{C. Tables moved from the main text}
\begin{table}[htbp]
\caption{Five-axis sweep on TUAB eval ($n{=}276$): AUROC of the naive and
canonicalized arms, mean over three MLP seeds. Stars on the canonicalized
value: canonicalized-vs-naive DeLong contrast on the seed-7 probe
surviving Benjamini--Hochberg at $q{=}0.05$ ($^{*}p<0.05$, $^{**}p<0.01$,
$^{***}p<0.001$). Identity rows omitted.}
\label{tab:s4-sweep}
\centering\scriptsize\setlength{\tabcolsep}{2pt}
\IfFileExists{../agg/table_multi_full.tex}{\input{../agg/table_multi_full.tex}}{\emph{pending}}
\end{table}

\begin{table}[htbp]
\caption{Failure-category counts on TUAB eval per backbone and arm
(seed-7 probes; the $12$ non-identity cells of A1--A3), with the number
of ground-truth normal recordings in parentheses.}
\label{tab:s5-failure}
\centering\tiny\setlength{\tabcolsep}{1.6pt}
\IfFileExists{../agg/table_failure_multi.tex}{\begin{tabular}{lcccccc}
\toprule
Category & \multicolumn{2}{c}{EEGPT} & \multicolumn{2}{c}{LaBraM} & \multicolumn{2}{c}{BIOT} \\
 & Naive & Canon. & Naive & Canon. & Naive & Canon. \\
\midrule
Robust-correct & 91 (0N) & 192 (89N) & 13 (5N) & 119 (29N) & 96 (0N) & 120 (24N) \\
Shift-recovered & 4 (0N) & 0 (0N) & 6 (6N) & 0 (0N) & 15 (0N) & 0 (0N) \\
Calibration collapse & 0 (0N) & 29 (19N) & 2 (2N) & 24 (11N) & 3 (0N) & 10 (0N) \\
Axis-specific & 7 (0N) & 3 (2N) & 8 (6N) & 0 (0N) & 2 (0N) & 0 (0N) \\
Shift-induced flip & 1 (0N) & 0 (0N) & 0 (0N) & 0 (0N) & 0 (0N) & 0 (0N) \\
Catastrophic shift & 125 (110N) & 0 (0N) & 199 (106N) & 79 (79N) & 129 (124N) & 100 (100N) \\
Robust-wrong & 48 (40N) & 52 (40N) & 48 (25N) & 54 (31N) & 31 (26N) & 46 (26N) \\
\bottomrule
\end{tabular}
}{\emph{pending}}
\end{table}

\begin{table}[htbp]
\caption{The five operators and the contract stage each naive arm
disables. Operators act on raw microvolts at native rate $f_s$.}
\label{tab:s6-operators}
\centering\scriptsize\setlength{\tabcolsep}{3pt}
\resizebox{\columnwidth}{!}{%
\begin{tabular}{@{}llll@{}}
\toprule
Axis & Operator $A_\theta$ on raw $\mu$V & Left-inverse in contract & Stage disabled (naive) \\
\midrule
A1 rate & anti-aliased resample $f_s\!\to\!f_n$ & resample $f_n\!\to\!250$\,Hz (C2) & no resample \\
A2 gain & $\times g$, all channels & per-recording robust z (C3) & fixed-scale z \\
A3 mains & $+$ mains $+$ harmonics $<f_n/2$, per-electrode amp./phase & notch at mains (C5) & notch off \\
A4 montage & drop channels / bipolar export & CAR over present channels (C4) & CAR over all 19 \\
A5 SNR & $+$ white noise per channel & none & none (arms identical) \\
\bottomrule
\end{tabular}}
\end{table}

\begin{table}[htbp]
\caption{Four contemporaneous approaches to heterogeneous EEG
acquisition. ``Fixes'' means a deterministic preprocessing contract;
``learns'' means invariance is trained into a model; ``measures''
means acquisition shift is simulated and recovery is measured.}
\label{tab:s7-compare}
\centering\scriptsize\setlength{\tabcolsep}{3pt}
\resizebox{\columnwidth}{!}{%
\begin{tabular}{@{}lllll@{}}
\toprule
Approach & Harmonization & Backbones & Shift simulated & Stage attribution \\
\midrule
BRIDGE-EEG~\cite{chowdhury2026bridge} & fixes (62-ch, 1--50 Hz, STFT) & own SE-ResNets & no & no \\
PRISM~\cite{prism2026} & learns (multi-source MAE) & own & no & no \\
GCN montage~\cite{gcnharmonization2026} & learns (electrode graph) & task CNNs & no & no \\
ESF (this work) & fixes (19-ch, 250 Hz, CAR, robust z) & three public, frozen & five axes, raw $\mu$V & per stage, per backbone \\
\bottomrule
\end{tabular}}
\end{table}

\section*{D. Extended results}
\subsection{Pinsker bound validation}
\label{sec:pinsker_emp}
For each measured (axis, severity, arm) we compute the AUROC gap to the
clean reference and two distances between per-recording logit
distributions: the class-conditional total variations of
Theorem~1 (which require labels and are available here) and
the marginal total variation (which a deployed system can observe
without labels). \PINSKERNUMBERS\ The implied ratio $C_{\mathrm{arm}} =
|\Delta\mathrm{AUROC}| / \mathrm{TV}(P_0, P_\theta)$ is reported per
arm; because $C_{\max}$ is the maximum of these ratios over the same
arms, the statement that the marginal bound with $C_{\max}$ holds on
every arm is a consistency check by construction, not an independent
validation. A deployed monitor that raises an alarm when the marginal
logit shift exceeds a threshold calibrated on such a sweep is a
heuristic without the guarantee of Theorem~1.


\subsection{Can a label-free monitor substitute for the contract?}
\label{sec:monitor}
A deployed system that does not apply the contract could still try to
\emph{detect} off-contract recordings and abstain. We test the two
label-free statistics available at inference: the across-group
variance of the recording's logits (the abstention statistic of
\S III-B of the main text) and the calibrated confidence $|p-0.5|$. For each
naive-arm cell with at least ten catastrophic recordings we ask
whether the statistic ranks the catastrophic recordings above the
rest, and what abstaining on it buys. It does not work. The logit
variance ranks catastrophic recordings at chance (median AUROC $0.54$
across $11$ cells, range $0.35$--$0.71$); abstaining on the
highest-variance $20\%$ of recordings removes a median $24\%$ of the
catastrophic ones, and abstaining on $50\%$ removes $58\%$, both
indistinguishable from random abstention, with no consistent AUROC
gain on the retained set. Confidence is worse than useless for this
purpose (median AUROC $0.57$, but inverted on the sampling-rate and
broadband cells), because the catastrophic recordings are by
definition confidently wrong. Uncorrected acquisition shift produces
failures that look, to the model, like ordinary confident decisions.
That is the strongest argument in this study for applying the
contract rather than detecting its absence.


\subsection{Reliability of the clean probes}
Figure~\ref{fig:s-reliability} gives the reliability diagram of the
Platt-scaled EEGPT seed-7 probe on the clean eval split.
\begin{figure}[htbp]
\centering
\includegraphics[width=0.8\columnwidth]{fig_reliability.pdf}
\caption{Reliability diagram, EEGPT seed-7 MLP probe, clean ESF eval
split ($n{=}276$); ECE $0.121$ (10 equal-width bins), per-bin Wilson
$95\%$ intervals.}
\label{fig:s-reliability}
\end{figure}

\section*{E. Provenance, compute, version note and future work}
\textit{Provenance.} Backbones are the public braindecode checkpoints
\texttt{eegpt-pretrained} (Hugging Face revision \texttt{e41cb3ae},
SHA-256 prefix \texttt{fb34c206}), \texttt{labram-pretrained}
(\texttt{0563b6c6}, \texttt{53b752ed}) and
\texttt{biot-pretrained-prest-16chs} (\texttt{f85518e9},
\texttt{b5892144}), loaded with braindecode~1.5.2 (commit
\texttt{619c7c01}), PyTorch~2.6.0 with CUDA~12.4, MNE~1.13.2,
SciPy~1.17.1 and scikit-learn~1.9.1 under Python~3.11.16. The
train/eval patient-ID intersection is $0$ ($2{,}076$ training and
$253$ evaluation patients). Feature extraction and probe training are
deterministic across machines: the EEGPT seed-$7$ probe reproduced its
validation AUROC to six decimals ($0.912151$) on a second machine.


\textit{Compute.} One NVIDIA A100 (40\,GB) and one A10: corpus
transfer $13$\,min, per-backbone train-split feature extraction
$15$--$36$\,min, four probes under $15$\,min, and the $46$-cell cache
$1$--$3$\,h per backbone, dominated by EDF reading and canonicalization
on CPU rather than by the backbone. Every number in the paper
regenerates from the released caches without a GPU.

% ===================== VERSION NOTE =============================

\subsection{Version note (v2), in full}


v2 corrects v1 in six respects, all found by the author on inspection
of the v1 code after submission.
(i)~\emph{Injection point.} v1 applied the acquisition operators after
canonicalization, contrary to the stated design. The v1 sampling-rate
axis was therefore a spectral time-warp, the v1 gain arm never composed
forward and inverse operators, and v1 mains amplitudes were in robust-SD
units. v2 applies all operators to the raw microvolt recording at native
rate and reruns all $46$ cells.
(ii)~\emph{Probe training.} The v1 probe was trained on a
patient-stratified split of the full TUAB corpus that overlapped the
evaluation split. The v2 probe is trained on the train split only.
(iii)~\emph{Theorem~1} is restated in class-conditional form;
the v1 step to marginal total variation was invalid (Appendix~A.4).
(iv)~Section D states that $C_{\max}$ is fitted to the
arms on which it is checked.
(v)~The conflict-of-interest statement is corrected.
(vi)~The v1 backbone export embedded an unseeded $19{\to}19$
channel-projection adapter that is no longer available; v2 initialises
it from seed $7$, freezes it, and releases it.
\emph{Withdrawn:} the v1 losses at $128$ and $512$\,Hz, the montage
calibration recovery, and the empirical support claimed for
Proposition~2. \emph{Changed:} baseline AUROC $0.908 \to
0.904$; sampling rate is recovered at every tested rate; the failure
taxonomy is reported per arm. \emph{Unchanged:} format, byline,
taxonomy definitions, Appendix~A.5.

% ===================== ACKNOWLEDGEMENTS =========================

\subsection{Future work}
\textbf{Future work.}
Each of the following is scoped as a separate follow-up:
(1)~an adjudicated, labelled multi-recorder study at the partner
clinics, with intake-side abstention keyed on logit magnitude rather
than calibrated probability, to close the residual of
the legacy-recorder section of the main text;
(2)~cross-corpus replication on TUEP, TUSZ and TUSL using the released
per-recording protocol;
(3)~measurement of the band-limit stage proposed above as ESF v1.1,
per backbone;
(4)~lightweight fine-tuning ablations (LoRA, last-block unfreeze) to
test whether the class asymmetry shrinks with mild adaptation or
persists as a structural property of the pretrained representation;
(5)~extension of the recoverability framework to composed shifts and
to ICA-decomposable physiological artifacts (EMG, eye blink), where
the operator is approximately rather than exactly invertible.


\bibliographystyle{IEEEtranN}
\bibliography{SPMB2026_Dammu_final_v4}

\end{document}
